\chapter{Introduzione}

\section{Scopo del Progetto}

L'\textbf{AI Document Assistant} è un'applicazione web full-stack progettata per consentire
agli utenti di interagire in linguaggio naturale con i propri documenti PDF attraverso un
sistema di \emph{Retrieval-Augmented Generation} (RAG).

Il sistema risolve un problema pratico comune: la difficoltà di estrarre rapidamente
informazioni specifiche da grandi volumi di documenti. Invece di scorrere manualmente
decine di pagine, l'utente può semplicemente porre una domanda e ricevere una risposta
precisa, contestualizzata e citata con i riferimenti di pagina.

\section{Caso d'Uso}

\textbf{Scenario tipico:} un professionista carica una raccolta di report, contratti o
articoli scientifici. Vuole sapere, ad esempio:

\begin{itemize}
  \item \emph{"Qual è la clausola relativa alla rescissione anticipata?"}
  \item \emph{"Quali metodologie statistiche sono state utilizzate nell'articolo?"}
  \item \emph{"Riassumi i punti chiave del capitolo sull'architettura."}
\end{itemize}

Il sistema recupera automaticamente i passaggi più rilevanti e genera una risposta
coerente, citando le pagine di origine.

\section{Obiettivi Tecnici}

\begin{enumerate}
  \item \textbf{Caricamento e indicizzazione}: elaborazione automatica del PDF in
        pipeline di ingestion (parsing $\rightarrow$ chunking $\rightarrow$ embedding
        $\rightarrow$ storage vettoriale).
  \item \textbf{Chatbot RAG}: risposta in linguaggio naturale basata esclusivamente
        sul contenuto del documento.
  \item \textbf{Ricerca semantica}: ricerca per significato, non per keyword.
  \item \textbf{Deployment containerizzato}: l'intero sistema avviabile con un
        singolo comando \texttt{docker compose up}.
  \item \textbf{Configurabilità}: supporto per provider LLM multipli (OpenAI e Ollama
        per uso locale/privato).
\end{enumerate}

\section{Tecnologie Principali}

\begin{table}[H]
  \centering
  \begin{tabularx}{\textwidth}{lX}
    \toprule
    \textbf{Componente} & \textbf{Tecnologia} \\
    \midrule
    Backend API       & FastAPI 0.115 + Uvicorn \\
    Orchestrazione LLM & LangChain 0.3 + LangGraph 0.2 \\
    Modello linguistico & OpenAI GPT-4o-mini / Ollama Llama 3.2 \\
    Embedding         & text-embedding-3-small (1536 dim) / nomic-embed-text \\
    Vector Database   & Qdrant (ricerca coseno) \\
    Database relazionale & PostgreSQL 16 + SQLAlchemy 2.0 (async) \\
    Parsing PDF       & pypdf 5 \\
    Containerizzazione & Docker Compose \\
    \bottomrule
  \end{tabularx}
  \caption{Stack tecnologico del progetto}
\end{table}
